The \code{BPatch\_loopTreeNode} class provides a tree interface to a
collection of instances of \code{BPatch\_basicBlockLoop}
contained in a \code{BPatch\_flowGraph}. The structure of the tree
follows the nesting relationship of the loops in a
function's flow graph. Each \code{BPatch\_­loopTreeNode} contains a
pointer to a loop (represented by
\code{BPatch\_basicBlockLoop}), and a set of sub-loops (represented
by other \code{BPatch\_loopTreeNode} objects). The root
\code{BPatch\_­loopTreeNode} instance has a null loop member since a
function may contain multiple outer loops. The outer
loops are contained in the root instance's vector of children.

\begin{lstlisting}[language=C]
  void foo() {
    for(int x=0; x<10; x++) {
      // ...
      for(int y = 0; y<10; y++) {
        // ...
        for(int z = 0; z<10; z++) {
          // ...
        }
      }
    }
    for(int i = 0; i<10; i++) {
      // ...
    }
  }
\end{lstlisting}

Each instance of \code{BPatch\_loopTreeNode} is given a name that
indicates its position in the hierarchy of loops. The name
of each root loop takes the form of \code{loop\_x}, where x is an
integer from 1 to n, where n is the number of outer loops in
the function. Each sub-loop has the name of its parent, followed by a
.y, where y is 1 to m, where m is the number of
sub-loops under the outer loop.

The \code{foo} function will have a root \code{BPatch\_loopTreeNode},
containing a \code{NULL} loop entry and two \code{BPatch\_loopTreeNode}
children representing the functions outer loops. These children would
have names \code{loop\_1} and \code{loop\_2}, respectively
representing the x and i loops. \code{loop\_2} has no children.
\code{loop\_1} has two child \code{BPatch\_loopTreeNode} objects, named
\code{loop\_1.1} and \code{loop\_1.2}, respectively representing the
y and z loops.

\begin{apient}
  BPatch_basicBlockLoop *loop
\end{apient}

\apidesc{
  A node in the tree that represents a single
  \code{BPatch\_basicBlockLoop} instance.
}

\begin{apient}
  std::vector<BPatch_loopTreeNode *> children()
\end{apient}

\apidesc{
  The tree nodes for the loops nested under this loop.
}

\begin{apient}
  const char *name()
\end{apient}

\apidesc{
  Return a name for this loop that indicates its position in the
  hierarchy of loops.
}

\begin{apient}
  bool getCallees(std::vector<BPatch_function *> &v, BPatch_addressSpace *p)
\end{apient}

\apidesc{
  This function fills the vector v with the list of functions that
  are called by this loop.
}

\begin{apient}
  const char *getCalleeName(unsigned int i)
\end{apient}

\apidesc{
  This function return the name of the ith function called in the loop\'s body.
}

\begin{apient}
  unsigned int numCallees()
\end{apient}

\apidesc{
  Returns the number of callees contained in this loop\'s body.
}

\begin{apient}
  BPatch_basicBlockLoop *findLoop(const char *name)
\end{apient}

\apidesc{
  Finds the loop object for the given canonical loop name.
}
